\section{Method: A Lens-Structured Context MDP}
\label{sec:method}

We formulate deliberation as a finite-horizon discounted MDP
$(S, A, P, R, \gamma)$ induced by a frozen language model: $S$ the set of
reasoning-forest states, $A$ the typed actions of
\cref{sec:method-actions}, $P(s' \mid s, a)$ the transition probability
(sampled, never enumerated, because it integrates over the frozen model's
continuations), $R(s, a, s')$ the reward of \cref{sec:method-reward}, and
$\gamma \in (0, 1]$ the discount factor. Episodes run for at most
$T_{\max}$ turns (the depth limit). This section
specifies the state (\cref{sec:method-state}), the action space and lens
semantics (\cref{sec:method-actions}), the controller network
(\cref{sec:method-controller}), and the reward (\cref{sec:method-reward}).
Training is deferred to \cref{sec:training}.

\subsection{State: an immutable reasoning forest}
\label{sec:method-state}

The state is an immutable forest of typed nodes: \emph{goals} (the problem
is the root; subgoals are optional decompositions), \emph{hypotheses}
(tentative claims attached to a goal), \emph{facts} (claims with provenance
\texttt{given}, \texttt{derived}, or \texttt{verified}), and
\emph{verifier-call records} (the input, output, and success flag of each
call to the symbolic oracle, retained as part of the state). Every action
returns a new state; nodes carry stable integer identifiers so actions can
reference them.

Two properties are load-bearing. First, the distinction between
\emph{verified facts} and \emph{tentative hypotheses} is what gives
\textsc{Verify} its meaning: without it the controller cannot distinguish
``I proved this'' from ``I think this.'' Predicates such as
\textit{solved} and \textit{verified} are \emph{derived} from the
structure; they are information the policy may read, never a control-flow
constraint. Second, decomposition is \emph{optional scaffolding}: the policy
may split a goal, work on some children, ignore others, abandon a branch, or
commit the root directly from a single verified hypothesis. The only hard
gate in the system is the verifier-vouched witness check at commit time.
Connectives on decomposition edges (conjunctive versus exhaustive-case) are
recorded as non-gating metadata; whether checked aggregation helps is an
ablation, deferred with the move from answer verification to proof
verification.

A worked forest for $|x-3| + |x+2| = 7$ by case analysis, exercising every
node kind and both closure polarities, is:

{\small
\begin{verbatim}
GOAL solve |x-3| + |x+2| = 7          (root; connective: exhaustive_case)
+- GOAL case x < -2:   -(x-3) - (x+2) = 7
|  +- FACT (1 - 7)/2 = -3  [derived]  -> VerifierCall(ok) -> [verified]
|  (committed: witness is the verified fact; x = -3 checks x < -2)
+- GOAL case -2 <= x <= 3:   (3-x) + (x+2) = 5, constant
|  +- HYPOTHESIS "5 = 7 has a solution"        [never verified]
|  (abandoned: the branch contributes nothing)
+- GOAL case x > 3:   (x-3) + (x+2) = 7
   +- FACT (7 + 1)/2 = 4   [derived]  -> VerifierCall(ok) -> [verified]
   (committed: x = 4 checks x > 3)
(root committed on the aggregate max(4, -3) = 4 when the largest solution
 is asked; positive closure twice, negative closure once)
\end{verbatim}
}

\noindent The same forest is pinned by an executable integration test in
the code accompanying this paper.

\subsection{Actions: verbs as lenses}
\label{sec:method-actions}

An action is $a = (\textit{verb}, \textit{target})$, where the target is a
node identifier. A verb is fully specified by five fields: a
state-conditioned legality mask $\mathrm{avail}(s,t)$; a JSON schema
constraining the LLM output; a scoped read $\mathrm{view}(s,t)$ (the lens
\emph{get}): the projection of $s$ rendered into the LLM context; a
localized write $\mathrm{apply}(s,t,o)$ (the lens \emph{put}): the update to
the target's neighborhood given LLM output $o$; and a reward, zero for
almost every verb. The transition samples
$o \sim \mathrm{LLM}(\cdot \mid \mathrm{view}(s,t) \mathbin{+\!\!+} \mathrm{suffix}_v)$
and sets $s' = \mathrm{apply}(s,t,o)$, where $\mathrm{suffix}_v$ is the
verb's fixed instruction block, the output fields, their constraints, and
a one-shot example for $v$, appended ($\mathbin{+\!\!+}$ is string
concatenation) after the rendered view: the prompt is the scoped reading
of the state followed by what to do with it. Because $P(s' \mid s, a)$
integrates over LLM continuations, the MDP is generative: transitions are
sampled, never enumerated.

The verb set is six core verbs plus one optional
(\cref{tab:verbs}), one per operation family that the prior-art vocabulary
collapses to, with the propose family split by the epistemic status of what
is written back (a tentative hypothesis versus a derived fact).
\textsc{Commit} is recursive but terminal only at the root: the structural
fact $\textit{target} = \textit{root}$ is what makes a commit terminal and
verifier-gated, so no separate final-answer verb exists. \textsc{Abandon}
is negative closure (a forward prune retaining history), kept as an
emittable verb because at inference budget $B{=}0$ no search controller
exists to prune for the policy. There is no \textsc{Summarize} verb:
compressing a closed subtree to its statement and verdict is a rendering
property of the lens.

\begin{table}[t]
\centering
\small
\begin{tabular}{@{}llll@{}}
\toprule
Verb & Family & Realized by & Reward \\
\midrule
\textsc{Split} & structure & LLM & none \\
\textsc{FormHypothesis} & propose & LLM & none \\
\textsc{ComputeSubexpr} & propose & LLM & none \\
\textsc{Verify} & assess (symbolic) & verifier & gated, capped (sole non-terminal channel) \\
\textsc{Commit} & resolve ($+$) & promotion & terminal verdict at root only \\
\textsc{Abandon} & resolve ($-$) & bookkeeping & none \\
\textsc{Critique} (optional) & assess (LLM) & LLM & none, by construction \\
\bottomrule
\end{tabular}
\caption{The verb set. Only \textsc{Verify} and the root \textsc{Commit}
ever touch reward, and both are verifier-decided.}
\label{tab:verbs}
\end{table}

The default view for a subgoal target is \emph{grounded-local}: the problem
statement, the ancestor path from root to target, and the target's own
subtree; sibling subgoals and unrelated branches are hidden, and closed
subtrees render as statement plus verdict. Because both the read and the
write are scoped to the target's neighborhood, acting on a subgoal is
identical to stepping a smaller MDP rooted there: locality, modularity, and
recursion are one property, derived from the action interface.

\paragraph{The rendered observation carries the structure explicitly.}
Four properties of the serialization proved load-bearing in the formative
campaign (\cref{sec:preliminary}) and are part of the method, not
implementation detail. (i) The forest is rendered \emph{as a tree}:
children are indented under their parent and every non-root node carries a
\texttt{parent=} tag, so the ownership relation between a goal and its
evidence is present in the observation rather than recoverable only from
node-creation order. (ii) An open goal with a discharging witness is
annotated \texttt{(evidence: node $N$ verified)}, naming the exact node a
\textsc{Commit} would use. (iii) A rejected (masked) action produces an
explicit rejection notice in the next controller input; without it the
observation after a masked action is near-identical to the one before it,
and a low-temperature controller repeats the same choice indefinitely.
(iv) The turn counter shows the horizon (\texttt{TURN 7/16}). The formative
evidence for (i) is stark: under a flat id-ordered rendering of the same
state, a controller at roughly $5\times$ the price gained nothing (1/20 versus 1/20 on
case-analysis problems), while restoring the tree structure at fixed
controller capability improved outcomes four- to five-fold. No controller,
at any capability, can act on structure the observation erases; because
the same serialization feeds the learned encoder, we expect (though the
trained controller does not yet exist to confirm it) the requirement to
carry over. Two structural guards accompany the
render: an exhaustive \textsc{Split} names its full case set in one call,
so a goal that already has children is not re-splittable, and a
\textsc{Compute} whose statement duplicates a fact already attached to the
same goal is masked rather than appended.

\subsection{The controller: factored heads over a pluggable encoder}
\label{sec:method-controller}

The learned heads value the \emph{global} state while the frozen LLM
generates from the \emph{local} view. Because the pointer head must score a
variable, growing set of targets, the encoder emits a pair: a global
readout $g$ and one vector $h_k$ per live node. The verb head maps an
adapter embedding of $g$ to masked logits over verbs; the pointer head
scores each candidate node against a query from $(g, \textit{verb})$,
realizing the factored policy $\pi(a \mid s)$, the probability of taking
action $a$ in state $s$, as $\pi(\textit{verb} \mid s)\,\pi(\textit{target}
\mid s, \textit{verb})$; the Q head scores
$Q(s, \textit{verb}, \textit{target})$, the expected discounted return of
taking that pair in $s$, per legal pair. Promoting the target into the action is what fixes credit
assignment: a value over verbs alone averages over every target the LLM
might pick, so the agent never learns \emph{which node to act on}.

The encoder is a pluggable axis with three rungs of structural fidelity:
(i) the frozen LLM's last-layer hidden states over a marker-serialized
state, reading per-node vectors at node-marker positions and the global
readout at a terminal marker, all from one forward pass (the v1 default;
causal visibility is the known caveat); (ii) a decoupled text-embedding
endpoint; (iii) a graph neural network whose node features are initialized
from text embeddings of each node's statement, with typed edges, removing
the causal limitation. Heads and losses are unchanged across the three.
\stub{encoder hyperparameters and adapter dimensions, once fixed}

\subsection{Reward: the verifier-anchor invariant}
\label{sec:method-reward}

The symbolic verifier is the only reward source. The first version is
terminal-only:
\[
r_T =
\begin{cases}
+1 & \text{verifier accepts the committed root answer,}\\
\;\;\,0 & \text{depth limit with no root commit,}\\
-1 & \text{verifier rejects the committed root answer,}
\end{cases}
\]
with intermediate reward zero, except for one constant housekeeping term:
a masked action (an illegal or failed operation, for example a rejected
duplicate \textsc{Compute} or an unusable generation) advances the turn
and receives a fixed $-0.01$. This constant carries no task information,
is not produced by the verifier or by any model, and cannot be gamed
upward; it exists so that a wasted turn is strictly worse than a legal
one. The anchor invariant is therefore precisely: no \emph{learned} or
\emph{content-dependent} reward enters anywhere, and every non-constant
reward is verifier-decided. If sparsity blocks progress, the only
sanctioned remedy is a gated, capped intermediate reward through
\textsc{Verify} (novel promotions only, deduplicated per target, cumulative
cap, tautologies rejected). Two design rules follow the strongest available
evidence: DeepSeek-R1's rejection of learned process-reward models on
reward-hacking grounds \citep{deepseekr1}, and CRITIC's finding that
self-critique without external grounding is unreliable \citep{critic}.
Aesthetic signals may enter only as search priors or pruning heuristics,
never as Bellman targets: a policy optimized against ``looks like good
reasoning'' learns to look like good reasoning at the detector's
resolution. The moment reward comes from a model's opinion rather than the
verifier's check, the system under study has changed.
